\section{Conclusions}\label{sec:conclusion}
This thesis studies boundary-focused constraints as a way to improve the segmentation accuracy reached with limited labels and a specified training policy. The LTN perspective supplies an interpretable description of the desired output properties; the implemented loss and its numerical reductions determine the actual training intervention.

The audited SwinUNETR comparison shows a positive validation-selected effect across all twelve tested fold--seed pairs with ten labelled volumes per fold under a maximum-75-epoch policy. Longer controls on fold 0 substantially reduce the observed advantage, making schedule dependence a central finding rather than an incidental limitation. The available evidence supports improved performance under restricted training and earlier optimization progress, while leaving the best attainable generalization performance unresolved.

The five-volume MedSAM3 pilot extends the investigation to a pretrained model. Adding bands and edge supervision during LoRA fine-tuning improves mean whole-hippocampus Dice by 1.252 percentage points under a fixed 200-update protocol. This is evidence that the approach can be useful in the tested adaptation setting. Its robustness across the complete planned replication remains to be established.

The resulting conclusion is conditional and practically meaningful: explicit boundary supervision can improve low-data segmentation within a limited training schedule, and its value must be evaluated jointly with the model, objective, data selection and optimization policy. The present results do not establish that logical notation alone confers an advantage, that additional pretraining data are unnecessary, or that the same gain will persist after longer optimization and independent evaluation.
